\documentclass[11pt]{article}
\usepackage[a4paper,margin=27mm]{geometry}
\usepackage{amsmath,amssymb,amsthm}
\usepackage{graphicx}
\usepackage{xcolor}
\usepackage{microtype}
\usepackage[colorlinks=true,linkcolor=blue!55!black,urlcolor=blue!55!black,citecolor=blue!55!black]{hyperref}

\newtheorem{theorem}{Theorem}
\newtheorem{remark}{Remark}
\newcommand{\Lip}{\operatorname{Lip}_0}
\newcommand{\OF}{\overline{\mathcal F}}
\newcommand{\Lop}{\mathcal L}
\newcommand{\A}{\mathcal A}
\newcommand{\B}{\mathcal B}

\title{An isometric structural no-other-examples theorem\\
for a minimal Lipschitz composition problem}
\author{Automated mathematics run \texttt{fa\_banach\_001}, lane 11}
\date{22 July 2026}

\begin{document}
\maketitle

\begin{center}
\fbox{\parbox{0.88\linewidth}{\textbf{Status: partial solution, likely valid.}
The unrestricted question remains open. Under either of two natural
hypotheses on the maximal hull, the result rules out both a new operator
class and a nonstandard ideal norm.}}
\end{center}

\section{Source question}

Nahuel Albarracin and Pablo Turco ask, on page 11 of
\emph{On the Lipschitz operator ideal $\Lip\circ\A\circ\Lip$}
\cite{AT}, whether there is any minimal Banach operator ideal $\A$ other than
the approximable ideal $\OF$ such that
\[
   (\Lip\circ\A^{\min}\circ\Lip)^{\min}
\]
is of composition type. Since the question concerns minimal $\A$, one may
write $\A^{\min}=\A$. Their Proposition 2.17 proves the necessary class condition
\[
   (\A^{\max})^{\mathrm{sur}}=\Lop . \tag{1}
\]

\begin{figure}[h]
  \centering
  \includegraphics[width=\linewidth]{figures/open_problem_crop.png}
  \caption{The exact source passage and its necessary condition, page 11 of
  arXiv:2211.14240.}
\end{figure}

\section{Definitions}

For a Banach operator ideal $\B$, its surjective hull is characterized by
\[
 T\in\B^{\mathrm{sur}}(E,F)
 \quad\Longleftrightarrow\quad
 TQ_E\in\B(\ell_1(B_E),F),
\]
where $Q_E:\ell_1(B_E)\to E$ is the canonical quotient. The ideal is
\emph{surjective} if $\B=\B^{\mathrm{sur}}$ isometrically. It is
\emph{symmetric} if $T\in\B$ exactly when $T^*\in\B$, with equal ideal norms,
and \emph{injective} if an isometric embedding $i:F\to G$ satisfies
$iT\in\B$ only if $T\in\B$. These are the standard
Banach-ideal conventions \cite{Pietsch,DF}.

\section{Idea of the proof}

Composition type carries more norm information than (1): for every
finite-dimensional $E$ it forces
\[
 \|I_{\mathcal F(E)}\|_{\B}=\|I_E\|_{\A}. \tag{2}
\]
For $E=\mathbb R$, the left side is the $\B$-norm of the identity of
$\mathcal F(\mathbb R)=L_1(\mathbb R)$, while the right side is one. Every
separable Banach space is a metric quotient of a norm-one complemented copy
of $\ell_1$ in this $L_1$ space. This first upgrades (1) to the isometric
equality $\B^{\mathrm{sur}}=\Lop$; maximality then removes the separability
restriction. The surjective case collapses immediately. In the symmetric and
injective case, adjoints of the canonical quotient maps give the canonical
isometric evaluation embeddings into $\ell_\infty$ spaces. Injectivity removes
these embeddings and forces every identity into $\B$ with norm one.

\section{Partial classification theorem}

\begin{theorem}
Let $\A$ be a minimal Banach operator ideal and put $\B=\A^{\max}$. Suppose
that either
\begin{enumerate}
  \item $\B$ is surjective, or
  \item $\B$ is symmetric and injective.
\end{enumerate}
Then, isometrically,
\[
 (\Lip\circ\A\circ\Lip)^{\min}
 \text{ is of composition type}
 \quad\Longleftrightarrow\quad
 \A=\OF .
\]
In particular, the source question has neither an additional operator class
nor a nonstandard ideal norm in either structural case.
\end{theorem}

\begin{proof}
Assume first that $(\Lip\circ\A\circ\Lip)^{\min}$ is of composition type.
We begin by retaining the norms in Propositions 2.14--2.15 of \cite{AT}.
Write $\mathcal S=\Lip\circ\A\circ\Lip$. For finite-dimensional $E$, the
Dirac map $\delta_E:E\to\mathcal F(E)$ factors strongly through $I_E$, so
\[
 \|\delta_E\|_{\mathcal S}\leq \|I_E\|_{\A}.
\]
The composition-type hypothesis and Proposition 2.14 give isometrically
$\mathcal S^{\max}=\B\circ\Lip$. Since the composition norm of a Lipschitz
map is the $\B$-norm of its linearization,
\[
 \|I_{\mathcal F(E)}\|_{\B}
 =\|\delta_E\|_{\B\circ\Lip}
 =\|\delta_E\|_{\mathcal S^{\max}}
 \leq \|I_E\|_{\A}. \tag{3}
\]
On finite-dimensional components, $\A$ and $\A^{\max}=\B$ agree
isometrically. The Godefroy--Kalton lifting theorem \cite{GK} supplies a
linear isometry $U_E:E\to\mathcal F(E)$ with
$\beta_EU_E=I_E$, where $\beta_E:\mathcal F(E)\to E$ has norm one. Hence
\[
 \|I_E\|_{\A}=\|I_E\|_{\B}
 \leq \|I_{\mathcal F(E)}\|_{\B}.
\]
Together with (3), this proves (2).

We next strengthen Proposition 2.17 of \cite{AT} to
\[
   \B^{\mathrm{sur}}=\Lop \quad\text{isometrically}. \tag{4}
\]
Indeed, put $E=\mathbb R$ in (2). The normalization of a Banach operator
ideal and the canonical isometry
$\mathcal F(\mathbb R)\cong L_1(\mathbb R)$ give
\[
 \|I_{L_1(\mathbb R)}\|_{\B}=1. \tag{5}
\]
Let $G$ be separable. Choose a sequence dense in the unit sphere of $G$, with
every tail dense. The associated norm-one map $q:\ell_1\to G$ is a metric
quotient: the usual successive-approximation expansion of a vector gives,
for every $\varepsilon>0$, a preimage $a$ with
$\|a\|_1\leq(1+\varepsilon)\|q(a)\|$. Also, disjoint unit intervals give an
isometric copy $U\ell_1\subset L_1(\mathbb R)$ and the map taking interval
averages is a norm-one projection $P$ onto it. Thus $qP:L_1(\mathbb R)\to G$
is a norm-one metric quotient. By (5), $qP\in\B$ with $\B$-norm one. The
quotient property of the surjective hull now gives
\[
 \|I_G\|_{\B^{\mathrm{sur}}}=1. \tag{6}
\]
Finally, maximal and surjective hulls commute isometrically
\cite{DF,Pietsch}. Because $\B$ is maximal,
$\B^{\mathrm{sur}}=(\B^{\mathrm{sur}})^{\max}$. For arbitrary $T:E\to F$,
every finite-dimensional/cofinite compression of $T$ has separable domain;
(6) and the ideal property make its $\B^{\mathrm{sur}}$-norm equal to its
operator norm. The maximal-hull formula therefore gives
$\|T\|_{\B^{\mathrm{sur}}}=\|T\|$, proving (4).

If $\B$ is surjective, then $\B=\B^{\mathrm{sur}}=\Lop$ isometrically.

Now suppose instead that $\B$ is symmetric and injective. Fix an arbitrary
Banach space $E$. By (4), $\|I_{E^*}\|_{\B^{\mathrm{sur}}}=1$. The
definition of the surjective-hull norm therefore gives
\[
 Q_{E^*}:\ell_1(B_{E^*})\longrightarrow E^*
 \quad\text{in }\B,\qquad \|Q_{E^*}\|_{\B}=1.
\]
By symmetry,
\[
Q_{E^*}^*:E^{**}\longrightarrow \ell_\infty(B_{E^*})
 \quad\text{belongs to }\B,\qquad \|Q_{E^*}^*\|_{\B}=1.
\]
Precompose with the canonical embedding
$J_E:E\to E^{**}$. The resulting operator is
\[
 j_E:=Q_{E^*}^*J_E,
 \qquad j_E(x)=(x^*(x))_{x^*\in B_{E^*}},
\]
the canonical isometric evaluation embedding of $E$ into
$\ell_\infty(B_{E^*})$. Hence $j_E\in\B$ with ideal norm one. Isometric
injectivity of $\B$, applied to $j_E\circ I_E$, yields
$\|I_E\|_{\B}=1$.

As $E$ was arbitrary, the ideal property applied to $T I_E$ gives
$\|T\|_{\B}\leq\|T\|$ for every bounded $T$. The reverse inequality is part
of the Banach-ideal axioms. Thus $\B=\Lop$ isometrically in the second case.

Since $\A$ is minimal, all equalities in
\[
  \A=(\A^{\max})^{\min}=\Lop^{\min}=\OF
\]
are isometric.
Conversely, for the standard approximable ideal $\A=\OF$, Proposition 2.16
of \cite{AT} gives
\[
  (\Lip\circ\OF\circ\Lip)^{\min}=\OF\circ\Lip,
\]
which is of composition type. This proves both directions.
\end{proof}

\section{Verification and limitations}

The new quantitative step is (2). It uses the isometric equalities in source
Proposition 2.14 and the Godefroy--Kalton linear lifting. The proof of (4)
then supplies the norm information not written explicitly in source
Proposition 2.17. A reviewer should check the standard isometric commutation
of maximal and surjective hulls and the quotient-hull norm convention.

The theorem does not settle maximal hulls that are nonsurjective and fail
symmetry or injectivity. The source's $L_1$-factorable test case lies outside
the surjective branch and is ruled out there by a separate geometric
argument. Consequently this packet is a solved structural subcase, not a full
answer to the unrestricted existence question.

\section{Bounded novelty check}

The local result and attempt indexes were searched for arXiv:2211.14240 and
the core composition-type terminology. A web search on 22 July 2026 used the
exact source sentence and combinations of ``minimal Banach operator ideal'',
``surjective hull'', ``symmetric'', and ``injective''. It found the source
\cite{AT} and the preceding Galois-connection paper \cite{TV}, but no later
resolution or statement of this structural theorem. Novelty confidence is
moderate-low because the proof is a short consequence of standard operator
ideal facts and may be folklore.

\section{Human-review recommendation}

Recommended verdict: \emph{likely valid partial result}. A reviewer should
first check the exact norm identity (2), then the metric-quotient and
maximal-hull passage proving (4). No computational verification is needed.

\begin{thebibliography}{9}
\bibitem{AT}
N. Albarracin and P. Turco,
\emph{On the Lipschitz operator ideal $\operatorname{Lip}_0\circ\mathcal A
\circ\operatorname{Lip}_0$}, arXiv:2211.14240, especially Propositions
2.16--2.17 and page 11 (2022; revised 2023).

\bibitem{TV}
P. Turco and R. Villafane,
\emph{Galois connection between Lipschitz and linear operator ideals and
minimal Lipschitz operator ideals}, J. Funct. Anal. 277 (2019), 434--451;
arXiv:1807.10786.

\bibitem{Pietsch}
A. Pietsch, \emph{Operator Ideals}, North-Holland, 1980.

\bibitem{DF}
A. Defant and K. Floret, \emph{Tensor Norms and Operator Ideals},
North-Holland, 1993.

\bibitem{GK}
G. Godefroy and N. Kalton, \emph{Lipschitz-free Banach spaces},
Studia Math. 159 (2003), 121--141.
\end{thebibliography}

\end{document}
